
\documentclass{article}
\usepackage{colortbl}
\usepackage{makecell}
\usepackage{multirow}
\usepackage{supertabular}

\begin{document}

\newcounter{utterance}

\centering \large Interaction Transcript for game `cladder', experiment `full\_v1.5\_default', episode 5930 with refineonly{-}llama{-}3.1{-}8b{-}instruct.
\vspace{24pt}

{ \footnotesize  \setcounter{utterance}{1}
\setlength{\tabcolsep}{0pt}
\begin{supertabular}{c@{$\;$}|p{.15\linewidth}@{}p{.15\linewidth}p{.15\linewidth}p{.15\linewidth}p{.15\linewidth}p{.15\linewidth}}
    \# & \multicolumn{2}{c}{Player} && \multicolumn{2}{c}{Game Master} \\
    \hline

    \theutterance \stepcounter{utterance}  
    & & & \multicolumn{4}{p{0.6\linewidth}}{
        \cellcolor[rgb]{0.9,0.9,0.9}{
            \makecell[{{p{\linewidth}}}]{
                \texttt{\tiny{[P1$\langle$GM]}}
                \texttt{You are an expert in causal inference. The following question is not a typical commonsense query, but rather a meticulously designed question created by a professor specializing in causal inference, intended to assess the students' mastery of the course content.} \\
\texttt{Imagine a self{-}contained, hypothetical world with only the following conditions, and without any unmentioned factors or causal relationships: Drinking coffee has a direct effect on wife and alarm clock. Wife has a direct effect on alarm clock. We know that drinking coffee causes alarm set by wife. drinking coffee or alarm set by wife causes ringing alarm. Would the alarm rings the next morning if drinking coffee instead of not drinking coffee?} \\
\texttt{Start your answer with 'Yes' or 'No', followed by additional reasoning or evidence to support your explanation.} \\
            }
        }
    }
     \\ \\

    \theutterance \stepcounter{utterance}  
    & & \multicolumn{4}{p{0.6\linewidth}}{
        \cellcolor[rgb]{0.95,0.95,0.95}{
            \makecell[{{p{\linewidth}}}]{
                \texttt{\tiny{[P1$\leftrightarrow$P]}}
                \texttt{{[}DRAFT{]}} \\
\texttt{No} \\
\\ 
\texttt{The given conditions imply a causal chain: } \\
\\ 
\texttt{1. Drinking coffee → Alarm set by wife (because drinking coffee causes alarm set by wife)} \\
\texttt{2. Alarm set by wife → Alarm rings (because drinking coffee or alarm set by wife causes ringing alarm)} \\
\\ 
\texttt{However, the direct effect of drinking coffee on alarm ringing is mediated by the wife's action of setting the alarm. If drinking coffee causes the wife to set the alarm, then not drinking coffee would likely result in the wife not setting the alarm. } \\
\\ 
\texttt{Therefore, without the direct effect of drinking coffee on alarm ringing, and with the wife's action being the primary cause of the alarm ringing, it is reasonable to conclude that the alarm would not ring the next morning if drinking coffee instead of not drinking coffee.} \\
            }
        }
    }
    & \\ \\

    \theutterance \stepcounter{utterance}  
    & \multicolumn{4}{p{0.6\linewidth}}{
        \cellcolor[rgb]{0.561,0.737,0.561}{
            \makecell[{{p{\linewidth}}}]{
                \texttt{\tiny{[P1$\rangle$GM]}}
                \texttt{No} \\
\\ 
\texttt{The given conditions imply a causal chain: } \\
\\ 
\texttt{1. Drinking coffee → Alarm set by wife (because drinking coffee causes alarm set by wife)} \\
\texttt{2. Alarm set by wife → Alarm rings (because drinking coffee or alarm set by wife causes ringing alarm)} \\
\\ 
\texttt{However, the direct effect of drinking coffee on alarm ringing is mediated by the wife's action of setting the alarm. If drinking coffee causes the wife to set the alarm, then not drinking coffee would likely result in the wife not setting the alarm. } \\
\\ 
\texttt{Therefore, without the direct effect of drinking coffee on alarm ringing, and with the wife's action being the primary cause of the alarm ringing, it is reasonable to conclude that the alarm would not ring the next morning if drinking coffee instead of not drinking coffee.} \\
            }
        }
    }
    & & \\ \\

    \theutterance \stepcounter{utterance}  
    & & & \multicolumn{2}{p{0.3\linewidth}}{
        \cellcolor[rgb]{0.9,0.9,0.9}{
            \makecell[{{p{\linewidth}}}]{
                \texttt{\tiny{[GM$|$GM]}}
                \texttt{No} \\
            }
        }
    }
    & & \\ \\

    \theutterance \stepcounter{utterance}  
    & & & \multicolumn{2}{p{0.3\linewidth}}{
        \cellcolor[rgb]{0.9,0.9,0.9}{
            \makecell[{{p{\linewidth}}}]{
                \texttt{\tiny{[GM$|$GM]}}
                \texttt{yes} \\
            }
        }
    }
    & & \\ \\

    \theutterance \stepcounter{utterance}  
    & & & \multicolumn{2}{p{0.3\linewidth}}{
        \cellcolor[rgb]{0.9,0.9,0.9}{
            \makecell[{{p{\linewidth}}}]{
                \texttt{\tiny{[GM$|$GM]}}
                \texttt{game\_result = LOSE} \\
            }
        }
    }
    & & \\ \\

\end{supertabular}
}

\end{document}
